\documentclass[11pt]{article}
\usepackage[margin=1in]{geometry}
\usepackage[T1]{fontenc}
\usepackage{lmodern}
\usepackage{microtype}
\usepackage{amsmath,amssymb,amsthm}
\usepackage[hidelinks]{hyperref}

\newtheorem*{theorem}{Theorem}
\newcommand{\wstar}{\mathrm{weak}^{*}}

\title{The partially bicyclic semigroup \(S_2\) has \((F_{*s})\)}
\author{Literature-implied answer for arXiv:1207.4531}
\date{11 August 2026}

\begin{document}
\maketitle

\paragraph{Status.}
\textbf{Literature-implied affirmative answer (full answer to Question 3).}
The implication below combines the presentation of \(S_2\) with a later
nonlinear Ryll--Nardzewski theorem.  It is not presented as a new theorem of
the run.

\paragraph{Source question.}
Lau and Zhang, \emph{Fixed point properties for semigroups of nonlinear
mappings and amenability}, arXiv:1207.4531, ask in Question~3 on printed
page~34:
\begin{quote}
Does the partially bicyclic semigroup
\[
 S_2=\langle e,a,b,c:ab=ac=e\rangle
\]
have the fixed-point property \((F_{*s})\)?
\end{quote}
Here \((F_{*s})\) says that every norm-nonexpansive representation of the
discrete semigroup on a nonempty norm-separable weak-star compact convex
subset \(Q\) of a dual Banach space, with weak-star continuous action, has a
common fixed point.  Their convention is \(T_{st}=T_s\circ T_t\).

\paragraph{Decisive supporting theorem.}
Wiśnicki, \emph{Around the nonlinear Ryll--Nardzewski theorem},
arXiv:1903.12123v2, proves in Corollary~3.3 on printed page~9 that if \(Q\) is
a norm-separable weak-star compact convex subset of a dual Banach space, then
every weak-star continuous, norm-nonexpansive, norm-distal semigroup action on
\(Q\) has a common fixed point \cite{Wisnicki}.  (The supporting paper writes
``separable weak-star compact'' after explicitly identifying the relevant
hypothesis as norm separability.)

\begin{theorem}
The partially bicyclic semigroup
\(S_2=\langle e,a,b,c:ab=ac=e\rangle\) has property \((F_{*s})\).
\end{theorem}

\begin{proof}
Consider an arbitrary representation covered by \((F_{*s})\), acting on
\(Q\), and put
\[
 A=T_a,\qquad B=T_b,\qquad C=T_c.
\]
The defining relations and the representation convention give
\[
 AB=I_Q=AC.
\]
For all \(x,y\in Q\), nonexpansiveness gives
\[
 \|x-y\|=\|ABx-ABy\|
 \leq \|Bx-By\|\leq \|x-y\|.
\]
Thus \(B\) is an isometric embedding of \(Q\) into itself.  The identical
argument with \(C\) shows that \(C\) is an isometric embedding.

Let \(R\) be the semigroup of self-maps of \(Q\) generated by \(B\) and
\(C\). Every member of \(R\) is a composition of isometric embeddings and is
therefore an isometric embedding.  Consequently the \(R\)-action is
norm-nonexpansive and norm-distal: for \(x\ne y\),
\[
 \|Tx-Ty\|=\|x-y\|>0\qquad(T\in R).
\]
Every member of \(R\) is weak-star continuous, because \(B\) and \(C\) are.
Wiśnicki's Corollary~3.3 therefore supplies \(q\in Q\) fixed by all of
\(R\), in particular
\[
 Bq=q=Cq.
\]
The relation \(AB=I_Q\) now yields
\[
 Aq=ABq=q.
\]
Thus \(q\) is fixed by the three generators \(a,b,c\), and hence by every
element of \(S_2\).  Since the original representation was arbitrary,
\(S_2\) has \((F_{*s})\).
\end{proof}

\paragraph{Provenance and scope.}
The supporting paper cites Lau--Zhang as background, but contains no occurrence
of ``partially bicyclic'', \(S_2\), or Question~3.  The answer is therefore an
agent-identified application rather than an explicit later resolution.  It
settles Question~3 in full, for every representation in the definition of
\((F_{*s})\); it does not settle the other open questions in Section~6 of the
source paper.

\paragraph{Search evidence.}
The four lightweight run indexes were searched for arXiv:1207.4531, the exact
semigroup presentation, ``partially bicyclic'', \((F_{*s})\), and
norm-separable weak-star fixed-point variants.  Exact-phrase web/arXiv searches
used the same terms.  No explicit later answer to Question~3 was found.  A
local search of arXiv:1903.12123 located Theorem~3.1 and Corollary~3.3 and
confirmed that the supporting paper cites the source but does not state this
application.

\paragraph{Human-review recommendation.}
Verify the order convention \(T_{st}=T_sT_t\) in the source definition and
the norm-separability hypothesis in Wiśnicki's Corollary~3.3.  Once those two
items are checked, the displayed two-line isometry reduction is exhaustive.

\begin{thebibliography}{9}
\bibitem{LauZhang}
A. T.-M. Lau and Y. Zhang,
\emph{Fixed point properties for semigroups of nonlinear mappings and
amenability}, J. Funct. Anal. 263 (2012), 2949--2977;
\href{https://arxiv.org/abs/1207.4531}{arXiv:1207.4531}.

\bibitem{Wisnicki}
A. Wiśnicki,
\emph{Around the nonlinear Ryll--Nardzewski theorem}, Math. Ann. 377
(2020), 267--279;
\href{https://arxiv.org/abs/1903.12123}{arXiv:1903.12123v2}.
\end{thebibliography}

\end{document}
